\Needspace{22\baselineskip}\vspace{6mm}
\pdfbookmark[1]{Praktische Anleitung zur Herstellung künstlicher Magnete.}{fischer-magnete}
\begin{center}\small Göttingische gelehrte Anzeigen.\quad 1832.\quad September 10.\end{center}

\textit{Praktische Anleitung zur vorteilhaften Herstellung und Zusammenfügung künstlicher Magnete, besonders der Hufeisen, geraden Stäbe, Kompassnadeln und anderer Nadeln, sowie die neueste Entdeckung, wie man ihnen die höchste Anziehungskraft verleiht, für Naturforscher, Ärzte, Seefahrer, Techniker und alle anderen Arten von Metallarbeitern, von \textup{\textsc{Friedrich Fischer}}, Lehrer und praktischer Techniker. Heilbronn.\ \textup{1833.} In der Class’schen Buchhandlung.} (58 S.\ in 8°.\ Mit zwei lithografierten Tafeln.)

Unter den vielfältigen Erscheinungen der magnetischen Kraft erregt besonders das Tragen beträchtlicher Lasten durch künstliche Magnete, deren Gewicht nur einen sehr kleinen Teil des Gewichts dieser Lasten ausmacht, die Bewun\SourceLeaf{598}derung der Liebhaber. Wissenschaftlich ist es dagegen nur von untergeordnetem Interesse und wird im reichhaltigsten neueren Werk über Physik kaum mit einigen Worten erwähnt. Die Gestaltung der für diesen Zweck am besten geeigneten künstlichen Magnete in Hufeisenform gilt daher weniger als Sache der Physiker. Sie liegt eher in den Händen von Personen, die damit ihren Lebensunterhalt verdienen und bisweilen angeblich neue und eigene Methoden unter dem Siegel der Verschwiegenheit zum Kauf anbieten. Ein neuerer Fall dieser Art hat den Verfasser der vorliegenden kleinen Schrift zu eigenen Versuchen veranlasst. Es gereicht ihm zur Ehre, dass er deren Ergebnisse ohne Vorbehalt und Geheimniskrämerei veröffentlicht. Der bei weitem größte Teil der Schrift ist künstlichen Magneten in Hufeisenform gewidmet. Ihre günstigsten Proportionen, die Auswahl und Behandlung des Stahls und die anzuwendenden Streichmethoden werden verständlich beschrieben. Liebhaber können sicher sein, dass sie sich durch Befolgen der Anweisungen jederzeit solche Magnete mit sehr großem Tragvermögen verschaffen können.

Der Rezensent würde sich auf diese Empfehlung der vorliegenden Schrift beschränken, wenn ihm nicht gerade die vom Verfasser verwendeten Streichmethoden Anlass zu einigen eigenen Bemerkungen gäben. Gemeint sind mit der im Titel erwähnten neuesten Entdeckung vermutlich diese Methoden, obwohl der Verfasser selbst einräumt, dass sie auch anderweitig bekannt sein mögen. Die Bemerkungen können als eine nicht unwichtige Ergänzung von \textsc{Coulomb}s Erfahrungen mit den allen Physikern wohlbekannten Methoden von \textsc{Knight}, \textsc{Duhamel}, \textsc{Michell}, \textsc{Canton} und \textsc{Aepinus} gelten.

Zur Erregung des Magnetismus in einem herzustellenden Hufeisenmagneten verwendet der Verfasser einen bereits vorhandenen Magneten derselben Form. Sein Verfahren besteht aus zwei nacheinander auszuführenden Arbeitsgängen. In beiden streichen stets beide Pole gleichzeitig. Im ersten folgt der eine Pol des Streichmagneten dem anderen auf dessen Weg. Im zweiten dagegen wird der eine Pol auf dem einen Arm und der andere auf dem anderen von der Krümmung zum Ende geführt. Die dafür erforderliche Anordnung der Pole braucht hier nicht eigens angegeben zu werden. Vor dem zweiten Arbeitsgang rät der Verfasser außerdem, den zu bestreichenden Magneten zu erwärmen und die Arbeit bis zur Abkühlung fortzusetzen.

\SourceLeaf{599}Man sieht leicht, dass der erste Arbeitsgang mit dem von \textsc{Michell} erfundenen Doppelstrich genau übereinstimmt. Der zweite entspricht dagegen im Wesentlichen \textsc{Duhamel}s Verfahren. Die getrennten geraden Stäbe (oder Bündel von Stäben), die \textsc{Duhamel} zum Streichen verwendet, bieten allerdings einige Vorteile für eine stärkere Erregung, die man bei Verwendung eines einzigen Hufeisenmagneten einbüßt (besonders weil man dann nicht in der Mitte der Krümmung beginnen kann). Da nun \textsc{Canton}s Methode bekanntlich lediglich die Methoden von \textsc{Michell} und \textsc{Duhamel} verbindet, ist das Verfahren des Verfassers im Wesentlichen das \textsc{Canton}sche mit den Änderungen, die der Einsatz eines hufeisenförmigen Streichmagneten von selbst mit sich bringt. Wirklich neu ist daher nur das vorherige Erwärmen, dessen Wirksamkeit allerdings wohl erst noch weiter bestätigt werden muss. Der Rezensent hat bei einigen eigenen Versuchen überhaupt keine besondere Wirkung davon festgestellt.

Besonders zu beachten ist hier nun, dass die Physiker dem Beispiel \textsc{Coulomb}s folgen und die Methode von \textsc{Canton} überhaupt nicht als Verbesserung gelten lassen. Nach dem Urteil jenes berühmten Physikers bestimmt stets allein die zuletzt angewandte Methode die Intensität des erregten Magnetismus. \textsc{Michell}s Streichart voranzustellen wäre daher ganz überflüssig. Andererseits geht aus den Äußerungen unseres Verfassers hervor, dass er die Verbindung seiner beiden Arbeitsgänge für wesentlich hält. Der Rezensent erkennt gerne an, dass diese Äußerungen, die von unprätentiöser Wahrheitsliebe zeugen, ihn zuerst dazu veranlasst haben, an der Allgemeingültigkeit des Grundsatzes zu zweifeln, auf dem \textsc{Coulomb}s Urteil beruht. Sehr viele Versuche mit besonderen, höchste Genauigkeit gewährenden Prüfmitteln, die an anderer Stelle beschrieben werden sollen, haben diesen Zweifel vollständig gerechtfertigt.

Bekanntlich ist die Verbesserung von \textsc{Michell}s Streichmethode, die wir \textsc{Aepinus} verdanken, am wirksamsten. Bei etwas dickeren Stahlstücken bleibt jede andere Methode, auch die \textsc{Duhamel}sche, weit hinter ihr zurück. \textsc{Coulomb}s Versuche haben dies außer Zweifel gestellt. Deshalb verwenden die Physiker zur stärkstmöglichen Magnetisierung solcher Stahlstücke ausschließlich die Methode von \textsc{Aepinus}. Bemerkenswert und nach den bisher ange\SourceLeaf{600}nommenen Voraussetzungen unerwartet ist daher das Ergebnis, das die erwähnten Versuche des Rezensenten übereinstimmend erbracht haben. Die nach der Methode von \textsc{Aepinus} so stark wie möglich magnetisierten Stahlstücke gewinnen immer noch beträchtlich an Kraft, wenn man sie \textit{nachher} wiederholt nach \textsc{Duhamel}s Verfahren streicht, obwohl dieses für sich allein nur eine beträchtlich geringere Kraft erzeugen kann als die Methode von \textsc{Aepinus}. Der Rezensent begnügt sich hier damit, diese Tatsache mitzuteilen, und unternimmt keinen Erklärungsversuch, der im Übrigen recht naheläge. Diese Erfahrungen wurden zwar unmittelbar nur bei der Magnetisierung gerader Stäbe gemacht. Es besteht aber kein Zweifel, dass die Verbindung der Methoden von \textsc{Aepinus} und \textsc{Duhamel} auch bei hufeisenförmigen Lamellen den stärkstmöglichen Magnetismus hervorrufen muss. Will man diese Verbindung bequem in unveränderter Form anwenden, sind allerdings einige besondere Vorkehrungen nötig. Wer sie nicht treffen möchte oder keine geeigneten geraden Stäbe zur Hand hat, erhält auch bei etwas dickeren Lamellen immer noch eine sehr große Stärke, wenn er nach der Anweisung des Verfassers einen hufeisenförmigen Streichmagneten verwendet. Die Stärke bleibt unter dem erreichbaren Höchstwert. Der Magnet lässt sich dabei sehr bequem handhaben.

Was der Verfasser über die Magnetisierung gerader Stäbe sagt, beschränkt sich auf die Handgriffe, die man ausführen muss, wenn man sie mit einem Hufeisenmagneten bestreichen will. Dadurch erhält man zwar eine große Stärke, doch die oben erwähnte Abfolge der Methoden von \textsc{Aepinus} und \textsc{Duhamel} ergibt eine noch größere und lässt auch an Bequemlichkeit nichts zu wünschen übrig.

Der Verfasser empfiehlt, magnetisierte gerade Stäbe zur Aufbewahrung mit den gleichnamigen Polen aufeinanderzulegen. Das ist völlig abzulehnen, wenn sie möglichst viel von ihrer Kraft behalten sollen. Am besten legt man sie paarweise in geringem Abstand so nebeneinander, dass ungleichnamige Pole einander benachbart sind, und legt Anker aus ganz weichem Eisen von geeigneter Länge daran.

